% s03-sec2.tex —— §2 函数域的情形（原文 p.12–17）
\section{函数域的情形}
\srcpage{12}

算子 $\prod_{v\notin S}L_v(s,\rho)$ 在整体域 $F$ 是有限域 $\kappa=\mathbb F_q$ 上几何连通曲线 $X$ 的有理函数域时，有一个非常漂亮的解释。

此时有限集 $S$ 是 $\kappa$ 上光滑射影曲线 $X$ 的一个闭子概形。设 $U$ 为其补集。

展开无穷乘积 $\prod_{v\notin S}L_v$，再按每个自然数 $d$ 把带因子 $q^{-ds}$ 的项归并，就得到形式级数的恒等式
\begin{equation}\tag{2.1}
\prod_{v\notin S}L_v(s,\rho)
=\sum_{d\in\mathbb N}q^{-ds}
\sum_{\sum_i\mathrm{div}_i\in U_d(\kappa)}\prod_i K^{\rho,(d_i)}_{v_i},
\end{equation}
其中 $U_d$ 表示 $U$ 的 $d$ 次对称幂，第二个和取遍 $U$ 中次数为 $d$ 的有效除子。算子
\begin{equation}\tag{2.2}
K_{\rho,d}=\sum_{\sum_i\mathrm{div}_i\in U_d(\kappa)}\prod_i K^{\rho,(d_i)}_{v_i}
\end{equation}
有几何解释。

我们不取 $F$ 上定义的群 $G$，而取 $X$ 上一个光滑群概形，其纤维连通、一般纤维为 $G$，且其在开集 $U$ 上的限制是约化的、拟分裂的，并在一个有限平展覆盖上分裂。为方便起见，我们仍用同一个字母 $G$ 记这个约化群概形。事实上，我们只是用几何语言重述了对集合 $S$ 的条件。设 $\mathcal O_x$ 是 $\mathcal O_X$ 在点 $x$ 的形式完备化。

\srcpage{13}
对 $x\notin U$，给定这个群概形就固定了 $G(F_v)$ 的一个紧子群 $G(\mathcal O_x)$，它不必是极大的。这个群概形的选取，本质上等价于选取阿代尔群 $G(\mathbb A)$ 的一个紧开子群 $G(\mathcal O_{\mathbb A})$。双陪集空间
\[
G(F)\backslash G(\mathbb A)/G(\mathcal O_{\mathbb A})
\]
可以用 $X$ 上的 $G$-torsor 来解释。分类 $X$ 上 $G$-torsor 的栈 $\Bun_G$ 是局部有限型的 Artin 代数栈。\footnote{由于若干原因，我们缺少 $G$-torsor 模栈上 $\ell$-进层的基础。一方面，当时缺少关于 Artin 代数栈上 $\ell$-进层的适当理论。幸运的是，Laszlo 与 Olsson 建立了这一理论（[LO]）。另一方面，所谓几何 Langlands 理论的专家们关心的是特殊尖形式的构造；对他们的目的而言，基础并非严格必需。由于我们关心迹公式，也就是同时关心所有自守形式，适当的基础对我们不可或缺。我们希望在后续文章中把它建立起来。目前可供参考的两个文献是 [LMB] 与 \url{http://www.math.upenn.edu/~kresch/teaching/stacks.html}。关于栈 $\Bun_G$ 本身，可参阅 Jochen Heinloth 的文章 [H]。} 其 $\kappa$-点集是一些双陪集空间的互不相交的并
\begin{equation}\tag{2.3}
\Bun_G(\kappa)=\bigsqcup_{\xi\in\ker^1(F,G)}G_\xi(F)\backslash G_\xi(\mathbb A)/G_\xi(\mathcal O_{\mathbb A}),
\end{equation}
其中 $G_\xi$ 是 $G$ 的一族形式，下标 $\xi$ 遍历 $H^1(F,G)$ 中对平展拓扑局部平凡的那些类构成的子集 $\ker^1(F,G)$，形式 $G_\xi$ 由与 $\xi$ 相伴的内扭构造。

Hecke 算子由修改（modification）栈上的一些 $\ell$-进反常层体现。记 $\Hecke$ 为四元组 $(x,E,E',\varphi)$ 的模栈，其中 $x$ 是 $X$ 的点，$E$ 与 $E'$ 是 $X$ 上的 $G$-torsor，$\varphi$ 是 $E$ 与 $E'$ 在 $X-\{x\}$ 上的限制之间的同构。换言之，$\varphi$ 是把 $E$ 在点 $x$ 处修改、结果为 $E'$ 的修改。至少起初，合理地限于修改位点 $\{x\}\subset U$。$\Hecke$ 在偶 $(x,E)$ 上的纤维此时是仿射 Grassmann 簇。它是一个带 $G(\mathcal O_x)$ 作用的 ind-概形。对偶群 ${}^LG$ 的任一有限维表示 $\rho:{}^LG\to\operatorname{GL}(V)$，都有 $\Hecke_U$ 上的一个复形 $A_\rho$，使得 $A_\rho$ 在 $\Hecke$ 于 $x\in U$ 与固定 $G$-torsor $E$ 上的纤维的限制，正是仿射 Grassmann 簇上与 $\rho$ 通过几何 Satake 等价（[MV]）对应的反常层。因此这个层的支集是有限维的，尽管纤维本身是无限维的。这些反常层 $K_\rho$ 组织成 $\Hecke$ 上的一个复形；经过适当的平移后它是反常层。$D^b(\Bun_G)$ 上的算子
\begin{equation}\tag{2.4}
\shF\longmapsto K_\rho(\shF)=\operatorname{pr}_{2,!}\bigl(\operatorname{pr}_1^*\shF\otimes K_\rho\bigr)
\end{equation}
就是 (2.2) 中 $K_{\rho,1}$ 在层论框架下的解释。这里 $\operatorname{pr}_1$ 与 $\operatorname{pr}_2$ 分别是 $(x,E,E',\varphi)$ 到分量 $E$ 与 $E'$ 的投影。

\srcpage{14}
借助 (2.3) 以及纤维上的 Frobenius 迹
\[
\sum_i(-1)^i\operatorname{tr}\bigl(\mathrm{Fr},H^{(i)}(F_x)\bigr),
\]
后者允许按 Grothendieck 的字典从一个层复形过渡到一个函数，经典的自守形式理论就变成了几何理论，只不过复数域被 $\bar{\mathbb Q}_\ell$ 的代数闭包代替（[Lm2]）。我们提醒，通常的 Hecke 算子有几何解释
\[
\shF\longmapsto\Hecke_\rho(\shF)
=\operatorname{pr}_{\Bun_G\times U,!}\bigl(\operatorname{pr}_1^*F\otimes K_\rho\bigr),
\]
其中 $\operatorname{pr}_{\Bun_G\times U}$ 是投影 $(x,E,E',\varphi)\mapsto(E',x)$。于是算子 (2.4) 就是把 $\Hecke_\rho(F)$ 沿 $U$ 积分。

为使 Arthur 参数
\begin{equation}\tag{2.5}
\sigma=\varphi\times\psi:\operatorname{SL}_2\times W_F\longrightarrow{}^LG
\end{equation}
定义一个 $\ell$-进层，必须把群 ${}^LG$ 解释为 $\bar{\mathbb Q}_\ell$ 上定义的群，而把 $W_F$ 解释为 Galois 群 $\operatorname{Gal}(F^{\mathrm{sep}}/F)$ 的子群，即其在 $\operatorname{Gal}(\bar F/F)$ 中的像是 $\mathbb F$ 的整数幂的那些元素构成的子群。首先，对 $\sigma$ 与 $\rho$ 附上 $\rho$ 的空间的一个分次结构，它由 $\operatorname{SL}_2$ 的对角环面作用定义，也就是由 $\sigma_j(1,-1)$ 的特征值 $\{m_j,m_j-2,\dots,-m_j\}$ 定义，其中 $\sigma_j$ 由 (1.2) 给出；然后附上一个分次局部系统，也就是 $W_F$ 的一个与分次结构相容的表示 $\psi'$。它由
\[
\psi':w\longmapsto\varphi
\begin{pmatrix}q^{n/2}&0\\0&q^{-n/2}\end{pmatrix}\psi(w)
\]
给出，其中 $w$ 在 $\operatorname{Gal}(F^{\mathrm{nr}}/F)$ 中的像是 $\mathbb F^n$。我们强调，我们假定 $\bar{\mathbb Q}_\ell$ 嵌入 $\mathbb C$，以便 $q^{n/2}$ 可以解释为 $\bar{\mathbb Q}_\ell$ 的元素。因此有一个古怪却熟悉的符号歧义。

\srcpage{15}
称层 $\shF$ 是关于以 Arthur 参数 (2.5) 为特征值的 Hecke 算子的\textbf{固有层}（propre），若存在同构\footnote{对经典理论更熟悉的读者，不妨停留在等式 (2.4) 上，说服自己它蕴含了与层相伴的函数的 Hecke 的经典定义，并且随后的关系蕴含：由迹与层相伴的函数，借助 [MV] 的同构，满足自 Hecke 以来已知的关系，只不过系数属于不同的域。}
\[
\Hecke_\rho(\shF)\simeq\shF\boxtimes L_\rho(\sigma),
\]
其中 $L_\rho(\sigma)$ 是把 $\rho\circ\sigma:W_F\times\operatorname{SL}_2\to{}^LG$ 复合后再限制到 $W_F$ 所得的分次局部系统。特别地，若 $\sigma$ 限制到 $\operatorname{SL}_2$ 是平凡的，则 $L_\rho(\sigma)$ 集中在 $0$ 次。此外，我们希望与 ${}^LG$ 的表示的张量积相容。若 $\shF$ 是关于以 Arthur 参数为特征值的 Hecke 算子的固有层，则
\[
K_\rho(\shF)\simeq\shF\otimes H^*_c(U,L_\rho(\sigma)).
\]
对每个 $d\in\mathbb N$，$d$ 次对称幂 $U_d$ 上的修改栈 $\Hecke_d$ 分类四元组 $(D,E,E',\varphi)$，其中 $D\in U_d$ 是 $U$ 中次数为 $d$ 的有效除子，$E$ 与 $E'$ 是 $X$ 上的 $G$-torsor，$\varphi$ 是 $E$ 与 $E'$ 在 $X-D$ 上的限制之间的同构。在一个固定 $G$-torsor $E\in\Bun_G$ 与一个固定的、约化支集为 $x_1+\cdots+x_r$、长度为 $r$ 的有效除子 $D=\sum_{i=1}^r d_ix_i$ 之上，$\Hecke_d$ 同构于 $r$ 个仿射 Grassmann 簇的积，每一个对应一个点 $x_i$。在 $\Hecke_d$ 之上，可以构造一个复形 $K_{\rho,d}$，使得在除子 $\sum d_ix_i$ 之上是 $\prod_{i=1}^rK^{\rho,(d_i)}_{x_i}$，其中 $K^{\rho,(d_i)}_{x_i}$ 是仿射 Grassmann 簇上与 $\rho$ 的 $d_i$ 次对称幂对应的反常层。这个反常层的存在并非理所当然，但将在后续文章 [FN] 中解释。注意，当 $G=\operatorname{GL}_n$ 且 $\rho$ 是标准表示时，它已由 Laumon 在 [Lm1] 中构造，一般构造与之相差不大。$D^b(\Bun_G)$ 上的算子
\[
\shF\longmapsto K_{\rho,d}(\shF)=\operatorname{pr}_{2,!}\bigl(\operatorname{pr}_1^*\shF\otimes K_{\rho,d}\bigr)
\]
就是 (2.2) 中 $K_{\rho,d}$ 在层论框架下的解释。这里 $\operatorname{pr}_1$ 与 $\operatorname{pr}_2$ 分别是 $(D,E,E',\varphi)$ 到分量 $E$ 与 $E'$ 的投影。

\srcpage{16}
\setcounter{theorem}{5}
\begin{lemma}
设 $\shF$ 是关于以 Arthur 参数 $\sigma$ 为特征值的 Hecke 算子的固有层，则
\begin{equation}\tag{2.7}
K_{\rho,d}(F)=F\otimes S^dH^*_c(U,L_\rho(\sigma)).
\end{equation}
\end{lemma}

\begin{proof}
由于这个引理只是作为我们策略的背景而给出，我们在此只满足于一个函数层面的论证，而把层等式 (2.7) 的证明留到文章 [FN]。记 $f$ 为按 Grothendieck 字典与层 $\shF$ 对应的 $G(F)\backslash G(\mathbb A)/G(\mathcal O_{\mathbb A})$ 上的函数。按定义，$K_{\rho,d}(f)$ 是 $U$ 中次数为 $d$ 的有效除子 $D$ 上的和
\[
\sum_{D\in U_d(\mathbb F_q)}K_{D,\rho}(f),
\]
其中附属于除子 $D=\sum d_ix_i$ 的算子 $K_{D,\rho}$ 是乘积 $\prod_iK^{\rho,(d_i)}_{x_i}$，其 Satake 变换是 ${}^LG$ 在 $\rho$ 的 $d_i$ 次对称幂 $\rho^{(d_i)}$ 上的表示的迹函数。由于 $\shF$ 是以 $\sigma$ 为特征值的固有层，$f$ 是每个算子 $K_{D,\rho}$ 的特征函数，特征值为
\[
\operatorname{tr}\Bigl(\mathrm{Fr}_q,\bigotimes_{i=1}^sS^{d_i}L_\rho(\sigma)_{x_i}\Bigr).
\]
于是 Grothendieck–Lefschetz 迹公式给出等式
\[
\sum_{D\in U_d(\mathbb F_q)}\operatorname{tr}\Bigl(\sigma_D,\bigotimes_{i=1}^sS^{d_i}L_\rho(\sigma)_{x_i}\Bigr)
=\operatorname{tr}\Bigl(\mathrm{Fr}_q,H^*_c(U_d\otimes_\kappa\bar\kappa,S^dL_\rho(\sigma)_{x_i})\Bigr),
\]
它正是该引理在函数层面的反映。
\end{proof}

当 $L_\rho(\sigma)$ 集中在 $0$ 次时，复形 $H^*_c(U\otimes_\kappa\bar\kappa,L_\rho(\sigma))$ 集中在 $0,1,2$ 次；若 $U$ 是仿射的，则只集中在 $1,2$ 次。若 $L_\rho(\sigma)$ 既没有常值子层也没有常值商层，则 $H^*_c(U\otimes_\kappa\bar\kappa,L_\rho(\sigma))$ 集中在 $1$ 次。对称幂 $S^dH^*_c(U,L_\rho(\sigma))$ 由以下公式计算：
\[
\bigoplus_{d_0+d_1+d_2=d}S^{d_0}H^0_c\otimes\wedge^{d_1}H^1_c\otimes S^{d_2}H^2_c[-d_1-2d_2],
\]
其中省略了紧跟在 $H^*_c$ 之后的显然的 $(U\otimes_\kappa\bar\kappa,L_\rho(\sigma))$。在函数层面，次数的平移几乎只是形式上的。若 $L_\rho(\sigma)$ 既没有常值子层也没有常值商层，则当 $d$ 大于 $\dim H^1_c(U\otimes_\kappa\bar\kappa,L_\rho(\sigma))$ 时 $S^dH^*_c(U,L_\rho(\sigma))=0$。

\srcpage{17}
这个维数可以用 $L_\rho(\sigma)$ 的野分歧，按 Grothendieck–Ogg–Shafarevich 公式 [R] 计算。
